\section{数学}

\subsection{组合数线性预处理}

\begin{lstlisting}[language=C++]
ll fac[N], inv[N], inv_fac[N];
void binom_init(int n) {
    fac[0] = fac[1] = inv[0] = inv[1] = inv_fac[0] = inv_fac[1] = 1;
    for(int i = 2; i <= n; i++) {
        fac[i] = fac[i - 1] * i % mod;
        inv[i] = (mod - mod / i) * inv[mod % i] % mod;
    }
    for(int i = 2; i <= n; i++) {
        inv_fac[i] = inv_fac[i - 1] * inv[i] % mod;
    }
}
int binom(int n, int m) {
    if (n < 0 || m < 0 || n < m) {
        return 0;
    }
    return fac[n] * inv_fac[m] % mod * inv_fac[n - m] % mod;
}
// by lyh
\end{lstlisting}

\subsection{线性筛}

\begin{lstlisting}[language=C++]
vector<int> primes, minp;
vector<int> phi;
void sieve(int n) {
    minp.assign(n + 1, 0);
    phi.assign(n + 1, 1);
    for (int i = 2; i <= n; i++) {
        if (!minp[i]) {
            primes.push_back(minp[i] = i);
            phi[i] = i - 1;
        }
        for (auto p : primes) {
            if (i * p > n) break;
            minp[i * p] = p;
            phi[i * p] = phi[i] * phi[p];
            if (p == minp[i]) {
                phi[i * p] = phi[i] * p;
                break;
            }
        }
    }
}
// by lyh
\end{lstlisting}

\subsection{扩展欧几里得}

$ax+by=c$ 有解 $\leftrightarrow (a,b)|c$，记 $x_0, y_0$ 为方程 $ax+by=(a,b)$ 的一组解，则原方程的解集为：
\[
    x = \dfrac{c}{g}x_0+k\dfrac{b}{g},\quad y = \dfrac{c}{g}y_0-k\dfrac{a}{g},\quad k\in \mathbb{Z}
\]

\begin{lstlisting}[language = C++]
ll exgcd(ll a, ll b, ll &x, ll &y){
    if(b == 0){
        x = 1; y = 0;
        return a;
    }
    ll g = exgcd(b, a%b, x, y);
    tie(x, y) = make_tuple(y, x-a/b*y);
    return g;
}
\end{lstlisting}

\subsection{高斯消元}

\begin{lstlisting}[language = C++]
double* solve(int n, double a[][N], double b[], int &state){
    state = 1;
    static double x[N];
    int rank = 0;
    memset(x, 0, sizeof(x));
    for(int i = 1; i <= n; i++){
        for(int j = 1; j <= n; j++){
            // 注意前一个未知数系数都为0的情况
            if(j < i && fabs(a[j][j]) > eps) continue;
            if(fabs(a[j][i]) < eps) continue;
            for(int k = 1; k <= n; k++) swap(a[i][k], a[j][k]);
            swap(b[i], b[j]);
            break;
        }
        if(fabs(a[i][i]) < eps) continue;
        ++rank;
        for(int j = 1; j <= n; j++){
            if(i == j || fabs(a[j][i]) < eps) continue;
            double r = a[j][i]/a[i][i];
            for(int k = i; k <= n; k++)
                a[j][k] -= a[i][k]*r;
            b[j] -= b[i]*r;
        }
    }
    for(int i = 1; i <= n; i++)
        if(fabs(a[i][i]) < eps && fabs(b[i]) > eps)
            {state = -1; return x;}
    if(rank < n) {state = 0; return x;}
    for(int i = 1; i <= n; i++)
        x[i] = b[i]/a[i][i];
    return x;
}
// 调用：double *x = solve(a, b, n, state);
\end{lstlisting}

\subsection{线性基}

A. 使用高斯消元构造基线性基，常数偏大，可处理第k小值。
\begin{lstlisting}[language = C++]
struct linear_base{
    using ll = long long;
    #define B 62
    ll a[B];
    linear_base(){
        memset(a, 0, sizeof(a));
    }
    void insert(ll x){
        for(int i = B-1; i >= 0; i--){
            if((x>>i&1) == 0) continue;
            if(a[i]) x ^= a[i];
            else{
                // 用 a[0...i - 1] 消去 x 的第 [0, i) 位上的 1
                for(int j = 0; j < i; j++)
                    if(x>>j&1) x ^= a[j];
                // 用 x 消去 a[i + 1...B-1] 的第 i 位上的 1
                for(int j = i+1; j < B; j++)
                    if(a[j]>>i&1) a[j] ^= x;
                a[i] = x;
                return;
            }
        }
    }
    ll get_max(){
        ll res = 0;
        for(int i = 0; i < B; i++)
            res ^= a[i];
        return res;
    }
};
\end{lstlisting}

B. 一般线性基，常数较小
\begin{lstlisting}[language = C++]
struct linear_base{
    using ll = long long;
    #define B 62
    ll a[B];
    linear_base(){
        memset(a, 0, sizeof(a));
    }
    void insert(ll x){
        for(int i = B-1; i >= 0; i--){
            if(x>>i&1){
                if(a[i]) x ^= a[i];
                else {a[i] = x; return;}
            }
        }
    }
    ll get_max(){
        ll res = 0;
        for(int i = B-1; i >= 0; i--)
            res = max(res, res^a[i]);
        return res;
    }
};
\end{lstlisting}

\subsection{快速卷积}

\subsubsection{FFT}

\begin{lstlisting}[language=C++]
using ld = double;
const ld PI = acos(-1);
struct Complex {
    ld x, y;
    Complex(ld x_ = 0, ld y_ = 0) : x(x_), y(y_) {}
    friend Complex operator+(const Complex &u, const Complex &v) {
        return {u.x + v.x, u.y + v.y};
    }
    friend Complex operator-(const Complex &u, const Complex &v) {
        return {u.x - v.x, u.y - v.y};
    }
    friend Complex operator*(const Complex &u, const Complex &v) {
        return {u.x * v.x - u.y * v.y, u.x * v.y + u.y * v.x};
    }
};

vector<int> rev;
void FFT(vector<Complex> &A, bool inverse) {
    int n = A.size();
    if ((int)rev.size() != n) {
        rev.resize(n);
        for (int i = 1; i < n; i++) {
            rev[i] = rev[i >> 1] >> 1 | (i & 1) << (__builtin_ctzll(n) - 1);
        }
    }
    for (int i = 0; i < n; i++) {
        if (i < rev[i]) {
            swap(A[i], A[rev[i]]);
        }
    }
    for (int len = 2; len <= n; len <<= 1) {
        Complex wn(cos(2 * PI / len), sin(2 * PI / len) * (inverse ? -1 : 1));
        for (int j = 0; j < n; j += len) {
            Complex w(1, 0);
            for (int k = j; k < j + len / 2; k++, w = w * wn) {
                Complex u = A[k], v = w * A[k + len / 2];
                A[k] = u + v;
                A[k + len / 2] = u - v;
            }
        }
    }
    if (inverse) {
        for (int i = 0; i < n; i++) {
            A[i].x = (int)(A[i].x / n + 0.5);
        }
    }
}
// by lyh
\end{lstlisting}

\subsubsection{NTT}

\begin{lstlisting}[language=C++]
constexpr int mod = 998244353;
int g0 = 3, g1;
vector<int> rev;
void NTT(vector<int> &A, bool inverse) {
    int n = A.size();
    if ((int)rev.size() != n) {
        rev.resize(n);
        for (int i = 1; i < n; i++) {
            rev[i] = rev[i >> 1] >> 1 | (i & 1) << (__builtin_ctzll(n) - 1);
        }
    }
    for (int i = 0; i < n; i++) {
        if (i < rev[i]) {
            swap(A[i], A[rev[i]]);
        }
    }
    for (int len = 2; len <= n; len <<= 1) {
        int gn = power(inverse ? g1 : g0, (mod - 1) / len);
        for (int j = 0; j < n; j += len) {
            int g = 1;
            for (int k = j; k < j + len / 2; k++, g = g * gn % mod) {
                int u = A[k], v = g * A[k + len / 2] % mod;
                A[k] = (u + v) % mod;
                A[k + len / 2] = (u - v + mod) % mod;
            }
        }
    }
    if (inverse) {
        int inv = power(n, mod - 2);
        for (int i = 0; i < n; i++) {
            A[i] = A[i] * inv % mod;
        }
    }
}

void solve() {
    g1 = power(g0, mod - 2);
    int n, m;
    cin >> n >> m;
    
    vector<int> a(n + 1), b(m + 1);
    for (int i = 0; i <= n; i++) {
        cin >> a[i];
    }
    for (int i = 0; i <= m; i++) {
        cin >> b[i];
    }
    int len = 1;
    while (len < n + m + 1) {
        len <<= 1;
    }
    a.resize(len);
    b.resize(len);
    NTT(a, 0);
    NTT(b, 0);
    vector<int> c(len);
    for (int i = 0; i < len; i++) {
        c[i] = a[i] * b[i] % mod;
    }
    NTT(c, 1);
    for (int i = 0; i <= n + m; i++) {
        cout << c[i] << " \n"[i == len - 1];
    }
}
// by lyh
\end{lstlisting}

\subsubsection{FWT}

\begin{lstlisting}[language=C++]
void Or(int n, int *a, int type) {
    for (int i = 1; i < n; i <<= 1) 
        for (int j = 0; j < n; j += i << 1) 
            for (int k = 0; k < i; k++) 
                (a[i + j + k] += a[j + k] * type + mod) %= mod;
}
void And(int n, int *a, int type) {
    for (int i = 1; i < n; i <<= 1) 
        for (int j = 0; j < n; j += i << 1) 
            for (int k = 0; k < i; k++) 
                (a[j + k] += a[i + j + k] * type + mod) %= mod;
}
void Xor(int n, int *a, int type) {
    for (int i = 1; i < n; i <<= 1) 
        for (int j = 0; j < n; j += i << 1) 
            for (int k = 0; k < i; k++) {
                (a[j + k] += a[i + j + k]) %= mod;
                (a[i + j + k] = a[j + k] - 2 * a[i + j + k] % mod + mod) %= mod;
                (a[j + k] *= type) %= mod;
                (a[i + j + k] *= type) %= mod;
            }
}
void FWT(int n, int *a, int f_type, int t_type) {
    if (f_type == 0) {
        if (t_type == 0) Or(n, a, 1);
        else Or(n, a, -1);
    } else if (f_type == 1) {
        if (t_type == 0) And(n, a, 1);
        else And(n, a, -1);
    } else {
        if (t_type == 0) Xor(n, a, 1);
        else Xor(n, a, power(2, mod - 2));
    }
}
// by lyh
\end{lstlisting}

\subsection{自适应辛普森}

辛普森公式：
\[
    \int_l^r f(x) \mathrm{d} x=\frac{(r-l)\left(f(l)+f(r)+4 f\left(\frac{l+r}{2}\right)\right)}{6}
\]
一般辛普森法：将区间 $[L, R]$ 分成 $2N$ 个小区间 $[x_{2i-1}, x_{2i}]$，其中 $x_i=l+i h, i=0 \ldots 2 n, h=\frac{r-l}{2 n}$.
\[
    \int_l^r f(x) d x \approx\left(f\left(x_0\right)+4 f\left(x_1\right)+2 f\left(x_2\right)+\ldots+4 f\left(x_{2 N-1}\right)+f\left(x_{2 N}\right)\right) \frac{h}{3}
\]

自适应辛普森时间复杂度：$O\left(\log \dfrac{1}{\epsilon}\right) \sim O\left(\dfrac{1}{\epsilon^{1/4}}\right)$
\begin{lstlisting}[language = C++]
double f(double x){
    // 给定函数
    return (c*x+d)/(a*x+b);
}

double simpson(double l, double r){
    return (r-l)*(f(l)+f(r)+f((l+r)/2)*4)/6;
}

double asr(double l, double r, double eps, double res, int step){
    double mid = (l+r)/2;
    double fl = simpson(l, mid), fr = simpson(mid, r);
    // 分段计算结果和不分段近似，且已完成强制次数，直接返回
    if(fabs(fl+fr-res)/15 < eps && step < 0)
        return fl+fr+(fl+fr-res)/15;
    return asr(l, mid, eps/2, fl, step-1)
        +asr(mid, r, eps/2, fr, step-1);
}

double integr(double l, double r){
    return asr(l, r, eps, simpson(l, r), 12);
    // 迭代次数限制可视情况改为6，注意卡常
}
\end{lstlisting}